% Part: first-order-logic
% Chapter: natural-deduction
% Section: derivations

\documentclass[../../../include/open-logic-section]{subfiles}

\begin{document}

\iftag{FOL}
      {\olfileid[pt]{fol}{ntd}{der}}
      {\olfileid[pt]{pl}{ntd}{der}}

\olsection{\usetoken{P}{derivation}}

\begin{explain}
Já dissemos o que é uma hipótese e demos as regras de
inferência.  !!^{derivation}s em dedução natural são geradas
indutivamente a partir disso: cada !!{derivation} ou é uma hipótese
por si só, ou consiste em uma, duas ou três !!{derivation}s seguidas
de uma inferência correta.
\end{explain}

\begin{defn}[!!^{derivation}]
\Article{derivation} \emph{!!{derivation}} de !!a{sentence}~$!A$ a partir de
hipóteses~$\Gamma$ é uma árvore finita de !!{sentence}s que satisfaz as
seguintes condições:
\begin{enumerate}
\item As !!{sentence}s no topo da árvore ou estão em $\Gamma$ ou
  são !!{discharged}s por uma inferência na árvore.
\item A !!{sentence} na base da árvore é~$!A$.
\item Toda !!{sentence} na árvore exceto a sentença~$!A$ na
  base é uma premissa de uma aplicação correta de uma regra de inferência
  cuja conclusão está diretamente abaixo daquela !!{sentence} na
  árvore.
\end{enumerate}
Dizemos então que $!A$ é a \emph{conclusão} da !!{derivation}
e $\Gamma$, suas hipóteses !!{undischarged}s.

Se existe !!a{derivation} de $!A$ a partir de~$\Gamma$, dizemos que $!A$ é
\emph{!!{derivable}} a partir de~$\Gamma$, ou, em símbolos: $\Gamma \Proves
!A$. Se há !!a{derivation} de~$!A$ em que toda hipótese é
!!{discharged}, escrevemos~$\Proves !A$.
\end{defn}

\begin{ex}
Toda hipótese por si só é !!a{derivation}. Assim, por exemplo, $!A$ por
si só é !!a{derivation}, e $!B$ por si só também. Podemos obter uma
nova !!{derivation} a partir delas aplicando, digamos, a regra $\Intro{\land}$,
\begin{prooftree}
\AxiomC{$!A$}
\AxiomC{$!B$}
\RightLabel{\Intro{\land}}
\BinaryInfC{$!A \land !B$}
\end{prooftree}
Essas regras devem ser gerais: podemos substituir o $!A$ e~$!B$
nela por quaisquer !!{sentence}s, por exemplo, por $!C$ e~$!D$. Então a
conclusão seria $!C \land !D$, e assim
\begin{prooftree}
\AxiomC{$!C$}
\AxiomC{$!D$}
\RightLabel{\Intro{\land}}
\BinaryInfC{$!C \land !D$}
\end{prooftree}
é uma !!{derivation} correta. É claro que também podemos trocar as
hipóteses, de modo que $!D$ desempenhe o papel de~$!A$ e $!C$ o
de~$!B$. Assim,
\begin{prooftree}
\AxiomC{$!D$}
\AxiomC{$!C$}
\RightLabel{\Intro{\land}}
\BinaryInfC{$!D \land !C$}
\end{prooftree}
também é uma !!{derivation} correta.

Podemos agora aplicar outra regra, digamos, $\Intro{\lif}$, que nos permite
concluir um condicional e nos permite descarregar qualquer hipótese
que seja idêntica ao antecedente desse condicional. Assim, ambas as
seguintes seriam !!{derivation}s corretas:
\begin{prooftree}
\AxiomC{$\Discharge{!C}{1}$}
\AxiomC{$!D$}
\RightLabel{\Intro{\land}}
\BinaryInfC{$!C \land !D$}
\DischargeRule{\Intro{\lif}}{1}
\UnaryInfC{$!C \lif (!C \land !D)$}
\DisplayProof\bottomAlignProof
\AxiomC{$!C$}
\AxiomC{$\Discharge{!D}{1}$}
\RightLabel{\Intro{\land}}
\BinaryInfC{$!C \land !D$}
\DischargeRule{\Intro{\lif}}{1}
\UnaryInfC{$!D \lif (!C \land !D)$}
\end{prooftree}
Elas mostram, respectivamente, que $!D \Proves !C \lif (!C \land !D)$ e
$!C \Proves !D \lif (!C \land !D)$.

Lembre-se de que descarregar hipóteses é uma permissão, não uma
exigência: não temos de descarregar as hipóteses. Em
particular, podemos aplicar uma regra mesmo que as hipóteses não
estejam presentes na !!{derivation}. Por exemplo, o seguinte é válido, embora
não haja hipótese~$!A$ a ser !!{discharged}:
\begin{prooftree}
  \AxiomC{$!B$}
  \DischargeRule{\Intro{\lif}}{1}
  \UnaryInfC{$!A \lif !B$}
\end{prooftree}
\end{ex}

\end{document}
